\documentclass[11pt]{article}
\usepackage[a4paper,margin=27mm]{geometry}
\usepackage{amsmath,amssymb,amsthm,hyperref}
\newtheorem{theorem}{Theorem}
\newtheorem{lemma}[theorem]{Lemma}
\newtheorem{corollary}[theorem]{Corollary}
\newcommand{\E}{\mathbb E}
\title{A logarithmic improvement in the lower bound for every die}
\author{Research note: endpoint odds, finite moments, and Laguerre tests}
\date{30 September 2026}
\begin{document}
\maketitle
\noindent\textbf{Status.} Complete proof, with three independent line-by-line mathematical
reviews completed. The endpoint-odds representation and the bad-prefix
estimate were developed in the parallel note
\texttt{../independent\_directions/endpoint\_poisson\_rigidity.tex}.
The finite-moment recurrence and the tapered Laguerre test below were
independently derived by collaborating agents. No formal verification is
claimed.

\begin{theorem}\label{main}
There are absolute constants $c>0$ and $m_0$ with the following property.
Let $p_1,\ldots,p_m$ be nondecreasing functions from
$\{1,\ldots,K\}$ to $[0,1]$, and let $w_q>0$ be probability weights.
Suppose
\begin{equation}\label{mixed}
 \E_w\prod_{j\in S}p_j=\frac1{|S|+1}
 \qquad(S\subseteq[m]).
\end{equation}
If $m\ge m_0$, then the last and first row weights satisfy
\[
 w_K\le\frac1{c m\log m},\qquad
 w_1\le\frac1{c m\log m}.
\]
In particular, an equal-weight representation has
\[
 K\ge c m\log m.
\]
\end{theorem}

\begin{corollary}\label{dice}
Every die in every permutation-fair collection of $n$ finite equiprobable
dice has at least $c n\log n$ faces, for an absolute $c>0$ and all
sufficiently large $n$. Thus the exceptional-die optimum satisfies
\[
 c n\log n\le g(n)\le Cn^2,
\]
where the upper bound is the independently established disjoint-bump
construction. Neither bound settles the minimum total size.
\end{corollary}
\begin{proof}
For a distinguished die, list its $K$ faces in increasing order and put
$p_j(q)=\Pr(X_j<X_0\mid X_0\text{ is face }q)$ for each other die $j$.
These functions are nondecreasing. Conditional independence and fairness
of every restriction give \eqref{mixed}, with $m=n-1$ and $w_q=1/K$.
\end{proof}

\section*{Endpoint odds and quantitative moments}
Put $A(q)=\prod_jp_j(q)$ and
\[
 \eta(q)=(m+1)w_q A(q).
\]
This is a probability measure by \eqref{mixed}. On rows with $A(q)>0$
write $r_j=(1-p_j)/p_j$ and $\lambda=\sum_jr_j$. Give a zero entry
infinite odds when defining bad rows; its $\eta$-mass is zero.

\begin{lemma}\label{moments}
Set $\varepsilon=m^{-1/2}$, and let $G$ consist of rows on which every
$r_j<\varepsilon$. For each integer $0\le k\le\sqrt m$, with $k<m$,
\[
 M_k:=\sum_{q\in G}\eta(q)\frac{\lambda(q)^k}{k!}
 \quad\text{satisfies}\quad
 |M_k-1|\le\frac{2k+1}{\sqrt m}.
\]
The measure restricted to $G$ need not have total mass exactly one.
\end{lemma}
\begin{proof}
Define $B_q(z)=\prod_j(p_j(q)+z(1-p_j(q)))$. Expanding
\eqref{mixed} gives the coefficient identities
\begin{align}
 (m+1)\E_w B_q(z)&=\sum_{k=0}^m z^k,\label{pgf}\\
 (m+1)\E_w\left[(1-p_j)\prod_{i\ne j}
                  (p_i+z(1-p_i))\right]
 &=\frac1m\sum_{k=0}^{m-1}(k+1)z^k.\label{derivative}
\end{align}
For example, the probability of a particular failure set of size $s$
is $\int_0^1 t^{m-s}(1-t)^s\,dt
=1/((m+1)\binom ms)$; summing over failure sets proves these identities.

The bad rows $D=G^c$ form an initial segment. If $D$ is nonempty,
choose a column $j_*$ with odds at least $\varepsilon$ at its last row.
Monotonicity gives $p_{j_*}\le(1-p_{j_*})/\varepsilon$ throughout $D$.
Consequently, as an inequality between nonnegative polynomial
coefficients,
\[
 (m+1)\E_w[B_q(z)\boldsymbol1_D]
 \preceq \frac{\varepsilon^{-1}+z}{m}
                 \sum_{k=0}^{m-1}(k+1)z^k.
\]
The same assertion is immediate when $D$ is empty. It includes rows
with zero entries without dividing by them.
Writing $e_k$ for the elementary symmetric polynomial in the odds,
\eqref{pgf} now implies, for $0\le k<m$,
\begin{equation}\label{ek}
 1-\frac{(k+1)/\varepsilon+k}{m}
 \le a_k:=\sum_{q\in G}\eta(q)e_k(r(q))\le1.
\end{equation}

On a good row, expand $\lambda^k$ as a sum over ordered $k$-tuples
of indices. The terms with distinct indices sum to $k!e_k$.
A union bound over a repeated pair of positions gives, for $k\ge2$,
\[
 0\le\lambda^k-k!e_k
 \le\binom{k}{2}\left(\sum_jr_j^2\right)\lambda^{k-2}
 \le\binom{k}{2}\varepsilon\lambda^{k-1}.
\]
Thus $0\le M_k-a_k\le((k-1)\varepsilon/2)M_{k-1}$; for $k=1$
there is equality $M_1=a_1$. Starting with $M_0\le1$, induction gives
$M_k\le1+k\varepsilon$ whenever $k\varepsilon\le1$: the preceding
moment is at most two, so $M_k\le1+(k-1)\varepsilon$.
The lower bound in \eqref{ek} and $M_k\ge a_k$ give
\[
 M_k\ge1-(k+1)m^{-1/2}-k/m\ge1-(2k+1)m^{-1/2}.
\]
This proves the lemma, including $k=0$.
\end{proof}

\section*{Two robust Laguerre tests}
Let $\rho(dx)=e^{-x}\,dx$ on $[0,\infty)$. Its normalized moments are
$\int x^k/k!\,d\rho=1$. Write
\[
 L_j(x)=\sum_{u=0}^j(-1)^u\binom ju\frac{x^u}{u!}.
\]
These Laguerre polynomials are orthonormal for $\rho$, and
\begin{equation}\label{recurrence}
 xL_j=(2j+1)L_j-(j+1)L_{j+1}-jL_{j-1},
\end{equation}
with the last term absent for $j=0$. Both facts follow from the displayed
formula: orthogonality also follows by $j$ integrations by parts in
$L_j(x)=e^x(j!)^{-1}(d/dx)^j(e^{-x}x^j)$, and the norm is one by
comparison with its leading coefficient.

For a polynomial $P(x)=\sum_s b_sx^s$, put
$\|P\|_*:=\sum_s|b_s|s!$. If a positive measure $\nu$ has
normalized-moment errors at most $\delta$ through degree $N$, then
\begin{equation}\label{error}
 \left|\int P\,d\nu-\int P\,d\rho\right|
 \le\delta\|P\|_*\qquad(\deg P\le N).
\end{equation}
Furthermore,
\begin{equation}\label{norms}
 \|L_iL_j\|_*\le3^{i+j},\qquad
 \|xP\|_*\le(\deg P+1)\|P\|_*.
\end{equation}
For the first bound, expand the two polynomials and use
$(u+v)!/(u!v!)\le2^{u+v}$.

\begin{lemma}\label{laguerre}
Suppose $d\ge2$ and $\nu$ is a finite positive measure on $[0,\infty)$
whose normalized moments through degree $2d-1$ differ from those of
$\rho$ by at most $\delta$. If
\[
 \delta d9^d<\frac16,
\]
then $\nu$ has positive mass in $[0,2/d)$. Moreover, every atom of
$\nu$ at a point $x\in[0,2/d]$ has mass at most
\[
 \frac{32}{d}+\frac{1024\delta9^d}{d^2}.
\]
\end{lemma}
\begin{proof}
Use the tapered polynomial
\[
 q_d(x)=\sum_{j=0}^{d-1}\left(1-\frac jd\right)L_j(x).
\]
Orthonormality and \eqref{recurrence} give
\begin{align*}
 \int q_d^2\,d\rho
 &=\frac{(d+1)(2d+1)}{6d},\\
 \int xq_d^2\,d\rho
 &=\sum_{j=1}^{d}j\left(\frac1d\right)^2
 =\frac{d+1}{2d}.
\end{align*}
Thus for $R_d(x)=(2/d-x)q_d(x)^2$,
\[
 \int R_d\,d\rho=\frac{(d+1)(d+2)}{6d^2}>\frac16.
\]
By \eqref{norms}, $\|q_d^2\|_*\le(\sum_{j<d}3^j)^2<9^d/4$,
and hence $\|R_d\|_*\le d9^d$ for $d\ge2$.
Equation \eqref{error} therefore gives $\int R_d\,d\nu>0$.
Since $R_d\le0$ on $[2/d,\infty)$, some positive mass lies below
$2/d$.

For the atom bound, fix $x\in[0,2/d]$ and set $r=\lfloor d/8\rfloor$.
For $0\le j\le r$, the coefficient formula yields
\[
 |L_j(x)-1|\le\sum_{u\ge1}\frac{(jx)^u}{u!}
 \le e^{1/4}-1<\frac12.
\]
Also $|L_j(x)|\le e^{1/4}<2$. Therefore
\[
 S_x:=\sum_{j=0}^r L_j(x)^2\ge\frac{r+1}{4}\ge\frac d{32}.
\]
The polynomial
\[
 P_x(y)=\frac1{S_x}\sum_{j=0}^r L_j(x)L_j(y)
\]
satisfies $P_x(x)=1$ and $\int P_x^2\,d\rho=1/S_x\le32/d$.
Its coefficient norm satisfies
\[
 \|P_x^2\|_*
 \le\frac1{S_x^2}\left(\sum_{j=0}^r |L_j(x)|3^j\right)^2
 \le\frac{1024}{d^2}9^{r+1}
 \le\frac{1024}{d^2}9^d.
\]
The mass of an atom at $x$ is at most $\int P_x^2\,d\nu$.
Apply \eqref{error} to finish the proof.
\end{proof}

\begin{proof}[Proof of Theorem~\ref{main}]
Take
\[
 d=\left\lfloor\frac{\log m}{8\log3}\right\rfloor,
 \qquad \nu=\sum_{q\in G}\eta(q)\delta_{\lambda(q)}.
\]
For all sufficiently large $m$, $d\ge2$ and $2d-1\le\sqrt m$.
Lemma~\ref{moments} gives the normalized-moment error
$\delta\le4d/\sqrt m$ through degree $2d-1$. Since $9^d\le m^{1/4}$,
\[
 \delta d9^d\le4d^2m^{-1/4}\longrightarrow0,
 \qquad
 \frac{1024\delta9^d}{d^2}
 \le\frac{4096m^{-1/4}}d=o(1/d).
\]
Lemma~\ref{laguerre} supplies a good row with $\lambda<2/d$.
All odds decrease with the row, so the last row is good and has
$\lambda(K)<2/d$. Its $\eta$-mass is bounded by the mass of the
corresponding atom of $\nu$, and therefore is at most $33/d$ for all
sufficiently large $m$.
Finally,
\[
 A(K)=\prod_j\frac1{1+r_j(K)}\ge e^{-\lambda(K)}\ge e^{-2/d}.
\]
Consequently
\[
 (m+1)w_K e^{-2/d}\le\eta(K)\le\frac{33}d.
\]
Since $d\asymp\log m$, this proves the last-row bound. Reversing row
order and replacing every $p_j$ by $1-p_j$ preserves \eqref{mixed}
(by expansion, or by uniform comparison patterns), so it also proves
the first-row bound. Equal weights now give $K\ge c m\log m$.
\end{proof}

\paragraph{Scope.}
The argument uses only the monotone local comparison tensor. It needs
neither rationality of its entries nor arithmetic divisibility. The
construction with $O(n^2)$ faces on one die shows that a substantially
stronger lower bound, if true, must remain compatible with that scale.
\end{document}
